\chapter{Introduzione}
\label{cap:introduzione}

\section{Contesto e obiettivo di ricerca}
\label{sec:intro-contesto}

\emph{One Health} parte da un'idea semplice: la salute di una popolazione non si può valutare separatamente dalle condizioni dell'ambiente in cui quella popolazione vive~\cite{ohhlep2022,salgado2020}. Tra i fattori ambientali, troviamo la qualità dell'aria ai primi posti: l'esposizione agli inquinanti atmosferici è riconosciuta da tempo come un determinante della salute, soprattutto respiratoria~\cite{hamanaka2025}. L'altro fattore preso in esame è il calore urbano: l'intensificarsi delle ondate di calore e l'effetto di isola di calore urbana (\emph{Urban Heat Island}, UHI) sono associati a un aumento della mortalità nelle città, in particolare tra gli anziani~\cite{iungman2023,huang2023}. 

I dati per studiare le correlazioni tra questi due fenomeni su scala locale non mancano, difatti, reti di monitoraggio ambientale, servizi di modellazione atmosferica e climatica e statistiche sanitarie ufficiali mettono a disposizione informazioni pubbliche e documentate. In questo contesto, il problema è la frammentazione delle informazioni, dovuta anche all' organizzazione dei dati in archivi separati e gestiti da enti diversi. A questo aggiungiamo anche la pubblicazione dei dati in formati eterogenei e riferiti a unità geografiche e temporali non coincidenti.
Quindi, mentre le misure ambientali sono puntuali (associate a una stazione) e con una risoluzione temporale elevata, gli indicatori sanitari sono aggregati per unità amministrative e per periodi più lunghi. Questo porta chi deve decidere a non avere la possibilità di leggerli insieme: ogni fonte è consultabile da sola, ma manca uno strumento che le metta in relazione sullo stesso territorio e nello stesso intervallo di tempo.

Per questo, il presente elaborato si rivolge a un destinatario preciso: i presidi medici e le aziende sanitarie, pubbliche e private, come supporto al controllo di gestione, tutti quegli attori quindi che potrebbero sfruttare l'aggregazione di questi dati per capire dove si concentra la domanda di prestazioni, come evolvono gli indicatori sanitari in rapporto a quelli ambientali e in quali aree ha più senso orientare risorse, prevenzione e investimenti. L'utenza è quindi professionale e orientata alla decisione, per questo il progetto non prevede un'interfaccia rivolta al singolo cittadino.

In questo senso, l'obiettivo del progetto è la realizzazione di una piattaforma software di supporto alle decisioni (\emph{Decision Support System}, DSS) che integri, a partire da soli dati pubblici, tre domini e li renda consultabili in modo unificato alla scala del distretto socio-sanitario: la \textbf{qualità dell'aria}, il \textbf{calore urbano} e gli \textbf{esiti di salute} (in particolare la mortalità). 

Mentre per gli inquinanti abbiamo fatto riferimento a CO, NO\textsubscript{2}, O\textsubscript{3}, PM\textsubscript{2.5} e SO\textsubscript{2} come principlai vettori di inquinamento, il calore è misurato tramite la temperatura dell'aria giornaliera per comune, affiancata da due fattori di mitigazione (copertura arborea e suolo impermeabile).

Il territorio di riferimento è il Salento, inteso in senso ampio come l'insieme delle province di Lecce, Taranto e Brindisi, suddiviso nei distretti socio-sanitari delle rispettive Aziende Sanitarie Locali~\cite{distretti_asl_lecce,distretti_asl_brindisi,distretti_asl_taranto}. Il distretto è stato scelto come unità di analisi perché corrisponde alla scala a cui le strutture sanitarie organizzano i servizi e pianificano le risorse, e perché offre un livello di aggregazione adeguato a trattare dati sanitari senza scendere a unità territoriali troppo piccole. Sul piano temporale l'analisi copre almeno tre anni, con l'obiettivo di avvicinarsi alla finestra 2016--2026 compatibilmente con la disponibilità effettiva dei dati. A questo rpoposito è bene precisare che sono stati utilizzati soltanto dati pubblici: nessuna rilevazione sul campo, nessun dataset fornito dall'azienda partner, nessuna credenziale di accesso a sistemi privati o ai gestionali delle ASL. Il vincolo è intenzionale, perché rende lo strumento replicabile e ogni elaborazione verificabile.
Il progetto, denominato OneHealth, è l'elaborato finale del Corso di Perfezionamento in Tecnologie Digitali per Ambienti e Sistemi Intelligenti dell'Università del Salento (IDA Lab), svolto in collaborazione con l'azienda partner Parsec 3,26 con sede a Lecce.

\section{Inquadramento e stato dell'arte}
\label{sec:intro-stato-arte}

La relazione tra ambiente e salute è già oggetto di diverse piattaforme istituzionali, esaminate nel Capitolo~\ref{cap:1}. A livello internazionale i riferimenti principali sono l'\emph{Environmental Public Health Tracking} dei CDC statunitensi~\cite{cdc_tracking} e l'Atlante europeo Ambiente e Salute dell'Agenzia Europea dell'Ambiente~\cite{eea_atlas}. In Italia il riferimento più vicino è l'Atlante Aria Salute (Ministero della Salute, ISS, ISPRA, Regioni)~\cite{atlante_aria_salute}.

Accanto alle piattaforme c'è un'ampia letteratura epidemiologica. Sul versante dell'aria, lo studio EpiAir2 associa l'inquinamento atmosferico ai ricoveri ospedalieri urgenti in 25 città italiane~\cite{scarinzi_epiair2_2013}. Sul versante del calore, uno studio su 12 città europee associa le alte temperature ai ricoveri per cause cardiovascolari e respiratorie~\cite{michelozzi_phewe_2009}, mentre studi su scala europea quantificano la quota di mortalità estiva attribuibile all'isola di calore urbana e il contributo di mitigazione della copertura arborea~\cite{iungman2023,huang2023}, mentre analisi su città mediterranee mostrano che l'effetto dipende fortemente dai fattori locali e va quindi studiato nella scala della singola area~\cite{cuerdo2023}. Per il territorio di riferimento, l'intensità dell'isola di calore di Lecce è stata misurata e pubblicata~\cite{pappaccogli_uhi_lecce_2024}, e costituisce un termine di paragone per i valori calcolati dal progetto. Questa letteratura conferma l'associazione tra fattori ambientali ed esiti di salute, ma richiama costantemente l'attenzione sui fattori confondenti, che rendono delicata ogni interpretazione.

Dall'analisi dello stato dell'arte non è emerso alcun precedente che combini insieme: 

- (i) un approccio interamente software basato su dati pubblici;

- (ii) l'integrazione congiunta di aria, calore ed esiti di salute; 

- (iii) un'unità di decisione a livello di distretto socio-sanitario, con utenza costituita da presidi e aziende sanitarie per il controllo di gestione; 

-(iv) la possibilità di interrogare i dati in linguaggio naturale tramite un modello linguistico eseguito in locale con tecniche di \emph{Retrieval-Augmented Generation} (RAG)~\cite{lewis_rag_2020}. 

Il precedente italiano più vicino, l'Atlante Aria Salute, opera su scala nazionale e regionale, non considera il calore urbano e non ha una componente conversazionale. Nessun dato UHI già pronto, inoltre, copre il Salento: i prodotti europei di riferimento (UrbClim) coprono un centinaio di città, in Puglia solo Bari, per cui l'isola di calore va ricostruita dai dati climatici pubblici.

\section{Contributo e soluzione proposta}
\label{sec:intro-contributo}

Il contributo del lavoro è una piattaforma organizzata in quattro componenti, descritte in dettaglio nei Capitoli~\ref{cap:2} e~\ref{cap:3}.

\paragraph{Acquisizione dei dati.}
I dati sono acquisiti da fonti pubbliche tramite chiamate dirette alle interfacce di programmazione degli enti, senza una copia persistente intermedia. Per la qualità dell'aria la fonte primaria è la rete di ARPA Puglia~\cite{arpa_puglia_api}, con OpenAQ, i dati dell'Agenzia Europea dell'Ambiente e il servizio modellato Copernicus/CAMS come fonti di riscontro e di copertura delle aree prive di stazioni~\cite{openaq,eea_aqds,cams_reanalisi,cams_previsioni}. Per il calore si usa la rianalisi climatica Copernicus ERA5-Land~\cite{era5_land}, con la rete meteo ARPA come validazione. Sul versante sanitario la fonte principale è ISTAT (decessi comunali)~\cite{istat_dati_mortalita}, con il bilancio demografico mensile ISTAT ed Eurostat come riscontro~\cite{istat_bilancio_mensile,eurostat_decessi_settimanali}.

\paragraph{Ingegneria dei dati.}
La seconda componente valida, pulisce e unifica i dati per zona, sovrapponendo nel tempo e nello spazio le serie ambientali e quelle sanitarie. L'unità di base è il comune di residenza, aggregato per distretto socio-sanitario. L'attendibilità si fonda sulla \textbf{ridondanza tra fonti indipendenti}: quando le fonti non concordano, la discordanza non viene nascosta ma tracciata e riportata come indicatore della qualità del dato. Dove mancano le stazioni, il dato ambientale disponibile è quello modellato e viene dichiarato come tale.

\paragraph{Mappa e dashboard.}
Una dashboard web eseguita in locale fornisce: indicatori sintetici, una mappa interattiva basata su OpenStreetMap~\cite{osm_copyright} con i confini reali dei distretti, una vista di dettaglio per distretto (che segnala anche la qualità del dato disponibile), serie temporali con selettore del periodo e una tabella comparativa. Quando un parametro non è coperto per un'area, l'interfaccia lo segnala esplicitamente invece di mostrare un valore non supportato dai dati.

\paragraph{Modulo di intelligenza artificiale.}
L'ultima componengte è un modulo basato su un modello linguistico (LLM) open-source e preaddestrato, eseguito in locale e senza alcun \emph{fine-tuning}, con architettura RAG vincolata ai soli dati della piattaforma. Risponde a domande in linguaggio naturale per un'inferenza di tipo esplorativo, indicando sempre le fonti su cui si basa. L'intera soluzione è pensata per essere ospitata in locale, senza servizi cloud, in risposta a un requisito di policy aziendale del committente Parsec 3,26: i dati elaborati e le interrogazioni rimangono sull'infrastruttura di chi usa lo strumento.

\section{Limiti e struttura della tesi}
\label{sec:intro-limiti}

Alcuni limiti vanno dichiarati fin dall'inizio. L'analisi è \textbf{esplorativa e non causale}: le associazioni rese visibili tra ambiente ed esiti di salute non vanno lette come rapporti di causa ed effetto, perché fattori noti come fumo, struttura per età e aria indoor non sono osservabili nelle fonti pubbliche usate. Le misure ambientali stimano inoltre l'esposizione \textbf{della popolazione}, non del singolo individuo: uno studio pilota del Politecnico di Milano, che confronta un sistema di sensori indossabili con la stima da stazione fissa, trova una dose inalata di PM\textsubscript{2.5} più alta del 22,3\%, con una differenza al limite della significatività statistica~\cite{bernasconi2025}: un compromesso accettato in cambio di scalabilità a costo zero. L'adozione del distretto come unità di presentazione, insieme al trattamento di soli dati aggregati, \textbf{mitiga il rischio di re-identificazione}. Infine, la \textbf{granularità e la copertura dei dati sanitari} variano da fonte a fonte e nel tempo, e i \textbf{tempi ridotti} hanno portato a privilegiare un prodotto completo dall'acquisizione dei dati all'interrogazione in linguaggio naturale, rispetto alla completezza delle singole funzionalità.

Il resto della tesi è organizzato come segue. Il Capitolo~\ref{cap:1} approfondisce il contesto e lo stato dell'arte. Il Capitolo~\ref{cap:2} è dedicato ai dati e al metodo: fonti, armonizzazione alla scala del distretto, pipeline e criteri di qualità del dato. Il Capitolo~\ref{cap:3} presenta l'architettura e l'implementazione della piattaforma, con attenzione alla dashboard, alla mappa e al modulo LLM/RAG. Le Conclusioni ( Capitolo~\ref{cap:conclusioni} ) riassumono il lavoro, ne discutono i limiti e indicano gli sviluppi futuri, tra cui la stima dell'intensità dell'isola di calore da dati satellitari.
